\documentclass[11pt,a4paper]{article}
\usepackage[utf8]{inputenc}
\usepackage{graphicx}
\usepackage[left=2.5cm,top=3cm,right=2.5cm,bottom=3cm,bindingoffset=0.5cm]{geometry}
\usepackage{AEDLogica, AEDEspecificacion, AEDTADs}
\usepackage{caratula}
\newlength{\miSangria}
\setlength{\miSangria}{2em}

% Asi pueden escribir nuevos comandos. 
% Este por ejemplo asegura q los nombres 
% que figuren con una tipografia diferenciada  
\newcommand{\Tipo}[1]{\mathsf{#1}}
% la sintaxis es \newcommand{\nombreDeLaMacro}[cantidadDeParametros]{Lo que va ser remplazado por el macro} 
\newcommand{\norm}[1]{\vert#1\vert}


\begin{document}

\section{IniciarPromocion}
	\vspace{1.0cm}

	\begin{proc}{IniciarPromocion}
		{
			\Inout AM: \Tipo{Airmiles},
			\In nombrePromo: \Tipo{\seq{char}},
            \In factorPromo: \Tipo{\R},
		}
		{}
		\comentario{Discutir el tercer elemento de la lista [factorPromo,0,0].
		NoPromo esta siempre en el sistema, no se puede pisar. Discutir si seria != noPromo o si directamente no tiene que pertenecer (podria existir y pisarse el valor?)}
		\\
		\requiere{AM = AM0 \land factorPromo \geq 1 \land nombrePromo \neq "NoPromo"}
		\aseguraLargo{
            AM.Usuarios = AM0.Usuarios \land
			\\
			AM.HistorialDeTransacciones = AM0.HistorialDeTransacciones \land
			\\
			AM.Promociones = SetKey(AM0.Promociones, nombrePromo, \lbrack factorPromo,0,-1 \rbrack)
		}
	\end{proc}

% Asegura{ AM.Usuarios = AM0.Usuarios  ∧ 
% 		   AM.HistorialDeTransacciones = AM0.HistorialDeTransacciones ∧
% 		   AM.Promociones = SetKey(AM0.Promociones, nombrePromo, [factorPromo,0,0] 		
% 		 }

\end{document}